\section{Demonstration}
\label{sec:demo}

Know-Surgery provides two interfaces for interactive exploration.

\paragraph{Command-Line Interface.}
Users run experiments with a single command, combining model, method, and dataset through Hydra overrides:

\begin{lstlisting}
# Edit: ROME on CounterFact with Qwen2.5-7B
python src/train.py mode=edit \
  trainer=edit/ROME model=Qwen2.5-7B-Instruct \
  experiment=edit/counterfact/default

# Unlearn: SimNPO on TOFU
python src/train.py mode=unlearn \
  trainer=SimNPO model=Llama-3.2-1B-Instruct \
  experiment=unlearn/tofu/default \
  task_name=my_simnpo_run

# Inject: LoRA on Alpaca
python src/train.py mode=inject \
  trainer=inject/LoRA model=Qwen2.5-7B-Instruct \
  experiment=inject/alpaca/default

# Evaluate a saved checkpoint
python src/eval.py \
  experiment=eval/tofu/default \
  model.model_args.pretrained_model_name_or_path=\
    saves/unlearn/my_simnpo_run
\end{lstlisting}

\paragraph{Web Interface.}
A Flask+Vue web UI (\texttt{new\_ui/}) allows users to select models, methods, and datasets from dropdown menus, configure hyperparameters, launch experiments, and compare results visually. The frontend is built with Vite + TypeScript and communicates with a Python Flask backend via REST APIs.

\begin{figure}[t]
\centering
\fbox{\parbox{0.95\columnwidth}{\centering\textit{[WebUI screenshot placeholder]}\\
\small Configuration panel + Results comparison view}}
\caption{Know-Surgery web interface for interactive experiment management.}
\label{fig:webui}
\end{figure}

\paragraph{Demo Video.}
A 2.5-minute screencast demonstrates the full workflow: selecting a method, configuring an experiment, running it, and analyzing results through both CLI and web interfaces.
